% =====================================================================
% METHODS INSERT (NEW SUBSECTION) — Epileptogenic-zone localization from the diffusion propagator
% Drafted 2026-07-07. Supports Results R3.1–R3.5 (epileptogenic markers).
%
% This is an INSERTION, not a find/replace: there is no existing epi-marker text in the
% Methods. Paste the \subsection block (below the horizontal rule) verbatim into the Methods,
% AFTER the cross-phase-comparison / anatomy subsections (ssec:methods_compare, sssec:methods_anatomy).
%
% BEFORE COMPILING
%   (1) Add ONE acronym to .agents/references/acronyms.tex (\gls{auc} is used in R3.1–R3.4 and
%       here, but is currently UNDEFINED):
%           \newacronym{auc}{AUC}{area under the receiver-operating-characteristic curve}
%       (glossary keys lrg, imcoh, soz, seeg already exist; roc is never used as \gls, skip it.)
%   (2) NO new math macros needed. Reuses: \lapl (L), \propag (K), \dm (\hat\rho), \partf (Z),
%       \Adjacency (W), \degm (D), \Tr, \diag, \abs, \ImCoh, \transp, \gls, \num, \qty.
%   (3) Cross-refs assume: ssec:methods_lrg defines the LRG propagator (\tau_{\min}=1/\lambda_{\max});
%       ssec:methods_compare is the cross-phase trace section. Fix the two \ref targets if the
%       Overleaf labels differ.
%
% WHAT EACH BLOCK SUPPORTS
%   sssec:methods_epi_affinity  -> R3.1 (relational, not nodal; one operator, second read-out)
%   sssec:methods_epi_gate      -> R3.1 (strength-matched gate; all-contacts + off-shaft designs)
%   sssec:methods_epi_detector  -> R3.2 (fused 6-band calibrated P(SOZ)); R3.3 depth baseline
%   sssec:methods_epi_seedfree  -> R3.3 (seed-free features), R3.4 (two-population switch), R3.5 (scope)
%
% GUARDRAILS honored (do not reintroduce on edit)
%   - NO results in Methods: no AUC/p-value/patient-count numbers. Those live in R3.1–R3.5.
%     Statistical PARAMETERS are kept (six bands, tau grid, R=300 surrogates, k=3 seeds, quartiles).
%   - Same e^{-tau L} operator as the trace, but read WITHIN one phase as a contact-to-contact
%     affinity from the raw kernel entries — NOT the distance/dendrogram/cophenetic/Spearman chain.
%     Stated once (block 1), never conflated (R3.1 header directive is emphatic on this).
%   - |ImCoh| is zero-lag-immune (Nolte 2004) -> matched-strength is the ONLY null in the text;
%     NO spatial/proximity null (feedback_no_spatial_proximity_null; the C5 spatial null is
%     belt-and-suspenders and deliberately omitted).
%   - Per-patient reporting, never pooled into a cohort scalar (never-pool rule; R3.4).
%   - NO repo/code technicalities (feedback_no_repo_or_code_technicalities_in_paper): no audit/
%     script/CSV names, no runtimes, no library notes in the LIVE text.
%   - Register: formal, present tense, first-person plural, em-dashes, \emph{}, \eqref/Section~\ref,
%     siunitx numbers, American English, one physical line per prose run.
%
% Provenance (AGENT-FACING ONLY — never enters the paper): audit_89/90 (relational co-diffusion
%   community + strength/spatial nulls), audit_101/102 (16-operator library, off-shaft LOSO + verify),
%   audit_132 (all-contacts headline, leave-one-contact-out, matched-strength), audit_117 (fused
%   LOPO detector, calibrated P(SOZ)), audit_104 (switch-not-blend), audit_118 (seed-free).
%   Math source: .agents/preprint/headlines/epi-marker-analysis.md §0.
% =====================================================================

% ---------------------------------------------------------------------
\subsection{Epileptogenic-zone localization from the diffusion propagator}
\label{ssec:methods_epi}

\subsubsection{Relational seizure-onset affinity from the within-phase propagator}
\label{sssec:methods_epi_affinity}
The cross-phase trace of Section~\ref{ssec:methods_compare} measures how far the task's reorganization of the diffusion hierarchy persists into rest; that hierarchy is generated by a propagator on the connectivity graph, and the same propagator, read directly within a single recording, supports a second read-out that localizes epileptogenic tissue. On each patient's resting (post-task) per-band graph \(\Adjacency = \abs{\ImCoh}\) --- symmetric, non-negative, zero-diagonal, one node per \gls{seeg} contact --- we form the combinatorial Laplacian \(\lapl = \degm - \Adjacency\), with \(\degm = \diag(s)\) and node strength \(s_i = \sum_{j}\Adjacency_{ij}\), and its heat-kernel propagator
\begin{equation}
    \propag(\tau) = e^{-\tau\lapl} = U\,\diag\!\big(e^{-\tau\lambda_{k}}\big)\,U^{\transp}, \qquad \dm(\tau) = \propag(\tau)/\partf, \quad \partf = \Tr\,e^{-\tau\lapl},
    \label{eq:methods_epi_kernel}
\end{equation}
with \((\lambda_{k},U)\) the eigendecomposition of \(\lapl\) (Section~\ref{ssec:methods_lrg}). Whereas the trace reads the normalized operator \(\dm(\tau)\) at the single fine scale \(\tau_{\min}=1/\lambda_{\max}\) and compares phases, the marker reads a single entry \(\big[\propag(\tau)\big]_{ij}\) --- the heat that reaches contact \(i\) from contact \(j\) in diffusion time \(\tau\) --- within one phase, so none of the communication-distance, dendrogram, cophenetic, or Spearman machinery of Section~\ref{ssec:methods_compare} enters here, and because the global partition function \(\partf\) is a scalar that leaves the ranking of contacts unchanged, the affinity is read from the unnormalized kernel. Scale is swept over a per-patient logarithmic grid of six diffusion times,
\begin{equation}
    \tau_{\ell} = \frac{10^{\,\ell/5}}{\lambda_{\max}}, \qquad \ell = 0,1,\dots,5,
    \label{eq:methods_epi_tau}
\end{equation}
so that \(\tau_{0}=1/\lambda_{\max}=\tau_{\min}\) is the fine scale of the trace and \(\tau_{5}=10/\lambda_{\max}\) the coarse scale at which the seizure marker is read. Let \(E\subset V\) be the set of seizure-onset contacts marked by the clinician on the implantation record, and let the affinity of a contact \(i\) to any target set \(T\) be the mean kernel weight to that set,
\begin{equation}
    \operatorname{aff}(i,T) = \frac{1}{\abs{T\setminus i}}\sum_{s\,\in\,T\setminus i}\big[\propag(\tau)\big]_{is}, \qquad \operatorname{aff}(i)\equiv\operatorname{aff}(i,E).
    \label{eq:methods_epi_aff}
\end{equation}
We read the propagator \emph{relationally}: a node-intrinsic reading --- each contact summarized by its own diffusion through the return probability \(\big[\propag(\tau)\big]_{ii}\), the diffusion-row entropy, or the commute distance --- is defined only as a comparison, and the seizure read-out is the affinity \(\operatorname{aff}(i)\) of a contact to the marked set. Because seizure tissue tends to be hyperconnected, node strength alone carries part of this signal; to keep the read-out orthogonal to hubness we residualize the affinity on strength, fitting \(\operatorname{aff}(i)\approx a\,s_{i}+b\) by least squares across contacts and retaining the residual
\begin{equation}
    \widetilde{\operatorname{aff}}(i) = \operatorname{aff}(i) - \big(a\,s_{i}+b\big).
    \label{eq:methods_epi_resid}
\end{equation}
Since the \gls{imcoh} substrate is immune to zero-lag volume conduction, this affinity cannot be a restatement of physical proximity, and a strength-matched null is the only control the connectivity read-out requires.

\subsubsection{Scoring and the strength-matched gate}
\label{sssec:methods_epi_gate}
Each contact score is turned into a marker by rank concordance with the labels: the \gls{auc} is the Mann--Whitney probability that a seizure contact outranks a healthy one,
\begin{equation}
    \mathrm{AUC} = \frac{\#\{\,\text{seizure}>\text{healthy}\,\} + \tfrac12\,\#\{\,\text{ties}\,\}}{n_{\text{seizure}}\;n_{\text{healthy}}}.
    \label{eq:methods_epi_auc}
\end{equation}
The primary evaluation is within patient and uses every marked contact: holding out one seizure contact at a time, we seed the affinity of \eqref{eq:methods_epi_aff} on all remaining seizure contacts and rank the held-out contact against every healthy contact by the strength-residualized score \(\widetilde{\operatorname{aff}}\) of \eqref{eq:methods_epi_resid}, admitting a patient only when at least four seizure and five healthy contacts are present. Significance is assessed against a strength-matched label null: \num{300} surrogate label sets, each matched to \(E\) in cardinality and in node-strength quartile, are passed through the identical leave-one-out marker to yield a per-patient upper-tail \(p\) and a cohort empirical \(p\) from the across-patient median score, and the cohort effect is a one-sided Wilcoxon signed-rank test of the per-patient \(\mathrm{AUC}-\tfrac12\). No spatial or proximity null is applied: it would be redundant on a substrate immune to zero-lag volume conduction, where matched strength is the mandatory control. As a conservative test of \emph{distant} discovery, we additionally evaluate off-shaft, partitioning contacts by electrode: for each electrode that carries seizure contacts we seed only on the seizure contacts of \emph{other} electrodes and score only contacts lying on electrodes that carry no seed, so that adjacency to a known contact cannot contribute and the read-out is tested where proximity is useless by construction.

\subsubsection{A fused, coordinate-free seizure-zone probability}
\label{sssec:methods_epi_detector}
To combine scales and rhythms into a single deployable read-out, we summarize each contact, relative to \(k=\num{3}\) known seizure seeds \(S\), by a vector of diffusion features: the heat-kernel affinity of \eqref{eq:methods_epi_aff} at fine, intermediate, and coarse \(\tau\), together with a small family of standard graph-diffusion kernels each read as a seed affinity --- personalized PageRank \((1-\alpha)\big(I-\alpha\,\degm^{-1}\Adjacency\big)^{-1}\), the Katz resolvent \((I-\beta\Adjacency)^{-1}\) at sub-critical attenuation \(\beta\), communicability \(e^{\Adjacency}\), the negative commute-time (effective resistance) \(-R\) with \(R_{ij}=\lapl^{+}_{ii}+\lapl^{+}_{jj}-2\lapl^{+}_{ij}\), and the diffusion distance between heat profiles --- and the seed-relative segregation \(\operatorname{aff}(i,S)-\operatorname{aff}(i,V\setminus S)\). No contact coordinates enter the feature set. The per-band features are concatenated across all six frequency bands, standardized within patient, and passed to an \(L_{2}\)-regularized logistic model trained leave-one-patient-out, producing a calibrated seizure-onset probability
\begin{equation}
    P(\mathrm{SOZ}\mid i) = \sigma\!\Big(b_{0} + \textstyle\sum_{f} b_{f}\,z_{f}(i)\Big),
    \label{eq:methods_epi_detector}
\end{equation}
with \(\sigma\) the logistic function and \(z_{f}\) the standardized features. The detector is scored by cross-patient \gls{auc}, by precision among its five highest-probability contacts and recall within its ten highest, and by the Brier score, all referred to the same strength-matched label shuffle as an empirical floor; we report the interpretable logistic model in preference to a higher-capacity nonlinear classifier, which is prone to overfitting at this cohort size.

\subsubsection{Seed-free features, the two-population switch, and validation}
\label{sssec:methods_epi_seedfree}
Two further reductions complete the analysis. First, a seed-free variant removes the dependence on known contacts, replacing the seed affinity with each contact's intrinsic diffusion signature --- the self-return \(\big[\propag(\tau)\big]_{ii}\), the localization of the slow Laplacian modes at the node (the inverse participation ratio of the low-\(\lambda\) eigenvectors), and its cophenetic persistence --- alongside a purely geometric baseline, the normalized along-shaft electrode depth; these are fused across bands and transferred leave-one-patient-out exactly as in \eqref{eq:methods_epi_detector}. Second, because a seizure network can present either as a co-diffusing community, in which the relational affinity is informative, or as a single high-strength hub, in which it is not, we select per patient between the two read-outs with a hard switch on the mean strength percentile \(r\) of the seeds,
\begin{equation}
    s_{\mathrm{switch}}(i) = \begin{cases} s_{i} & r \ge 0.70 \quad(\text{hub regime}),\\[2pt] \widetilde{\operatorname{aff}}(i) & r < 0.70 \quad(\text{community regime}), \end{cases}
    \label{eq:methods_epi_switch}
\end{equation}
rather than a continuous blend, which would average a working detector against a known failure mode. Performance is accordingly reported per patient and never pooled into a single cohort number. All read-outs are validated against the clinician-marked seizure-onset labels rather than against resection margins or surgical outcome, so the read-out is characterized as a connectivity marker that co-locates with the clinical seizure-onset zone, its precision interpreted as a triage shortlist over the implanted contacts.
